\subsection{EQUIVARIANT EMBEDDING THEOREM}

In this number, we assume that \(r \neq \omega\) .

Lemma 3. Let G be a compact topological group operating continuously on a topological space X; let A be a subset of X, stable under G, and W a neighbourhood of A. Then, there exists an open neighbourhood V of A stable under G and contained in W.

Put \(\mathrm{F} = \mathrm{X} - \mathrm{W}\) and \(\mathrm{V} = \mathrm{X} - \mathrm{GF}\) . Then \(\mathrm{V}\) is open (General Topology, Chap. III, §4, no. 1, Cor. 1 of Prop. 1), stable under \(\mathrm{G}\) , and \(\mathrm{A} \subset \mathrm{V} \subset \mathrm{W}\) .

THEOREM 1. Let G be a compact Lie group, \((g,x)\mapsto gx\) a law of left operation of class \(C^{r}\) of G on X, and A a compact subset of X. There exists an analytic linear representation \(\rho\) of G on a finite dimensional vector space E, a morphism \(\varphi:X\to E\) of class \(C^{r}\) , compatible with the operations of G, and an open neighbourhood U of A, stable under G, such that the restriction of \(\varphi\) to U is an embedding.

Replacing A by the compact set GA, we are reduced to the case in which A is stable under G.

Let \(E_0\) be a finite dimensional vector space such that there exists an element of \(\mathcal{C}^r(\mathrm{X};\mathrm{E}_0)\) that is an embedding in the neighbourhood of A (no. 1, Prop. 4); the set \(\mathcal{P}\) of morphisms having this property is a non-empty open subset of \(\mathcal{C}^r(\mathrm{X};\mathrm{E}_0)\) (no. 1, Prop. 3). Consider the continuous linear representation of the compact group G on the space \(\mathcal{C}^r(\mathrm{X};\mathrm{E}_0)\) (§6, no. 4, Lemma 4). By the Peter-Weyl theorem (Spectral Theory, in preparation), the union of the finite dimensional subspaces stable under G is dense in \(\mathcal{C}^r(\mathrm{X};\mathrm{E}_0)\) ; hence, there exists an element \(\varphi_0\) of \(\mathcal{P}\) such that the maps \(x \mapsto \varphi_0(gx)\) , for \(g \in G\) , generate a finite dimensional vector subspace \(E_1\) of \(\mathcal{C}^r(\mathrm{X};\mathrm{E}_0)\) , which is clearly stable under the operation of G.

Take E to be the space \(\operatorname{Hom}_{\mathbf{R}}(\mathrm{E}_{1},\mathrm{E}_{0})\) , \(\rho\) to be the representation of G on E induced by the action on \(E_{1}\) , and \(\varphi:X\to E\) to be the map that associates to \(x\in X\) the linear map \(\psi\mapsto\psi(x)\) from \(E_{1}\) to \(E_{0}\) . This is a morphism of class \(C^{r}\) ; for \(x\in X\) , \(g\in G\) , \(\psi\in E_{1}\) , we have (denoting by \(\tau(g)\) the automorphism \(x\mapsto gx\) of X):

\[
\varphi (g x) (\psi) - \psi (g x) = \varphi (x) (\psi \circ \tau (g)) = (\rho (g) \varphi (x)) (\psi).
\]

Let \(\alpha : \mathrm{Hom}_{\mathbf{R}}(\mathrm{E}_1, \mathrm{E}_0) \to \mathrm{E}_0\) be the linear map \(u \mapsto u(\varphi_0)\) ; we have \(\alpha \circ \varphi = \varphi_0\) , so \(\varphi\) is an embedding in the neighbourhood of A because \(\varphi_0\) is one. Hence, there exists an open neighbourhood U of A such that the restriction of \(\varphi\) to U is an embedding; we can choose U stable under G by Lemma 3, hence the theorem.

COROLLARY 1. Assume that X is compact. There exists an analytic linear representation \(\rho\) of G on a finite dimensional vector space E and an embedding \(\varphi: X \to E\) such that \(\varphi(gx) = \rho(g)\varphi(x)\) for \(g \in G, x \in X\) .

COROLLARY 2. Let H be a closed subgroup of G. There exists an analytic linear representation of G on a finite dimensional vector space E and a point \(v \in E\) with fixer H.

Apply Cor. 1 to the canonical operation of G on the compact manifold G/H. This gives an analytic linear representation \(\rho: \mathrm{G} \to \mathbf{GL}(\mathrm{E})\) and an embedding \(\varphi: \mathrm{G} / \mathrm{H} \to \mathrm{E}\) such that \(\varphi(gx) = \rho(g)\varphi(x)\) , \(g \in \mathrm{G}, x \in \mathrm{G} / \mathrm{H}\) . Let \(\bar{e} \in \mathrm{G} / \mathrm{H}\) be the class of \(e \in \mathrm{G}\) , and \(v = \varphi(\bar{e})\) its image. For all \(g \in \mathrm{G}\) , we have \(\rho(g)v = v \Longleftrightarrow \varphi(g\bar{e}) = \varphi(\bar{e}) \Longleftrightarrow g\bar{e} = \bar{e} \Longleftrightarrow g \in \mathrm{H}\) .

COROLLARY 3. Assume that X is paracompact. There exists a real Hilbert space E, a continuous unitary representation \({}^{10}\) \(\rho\) of G on E and an embedding \(\varphi: X \to E\) of class \(C^{r}\) such that \(\varphi(gx) = \rho(g)\varphi(x)\) for all \(g \in G\) and all \(x \in X\) .

The space X/G is locally compact (General Topology, Chap. III, §4, no. 5, Prop. 11). Its connected components are the images of the connected components of X, which are countable at infinity (General Topology, Chap. I, §9, no. 10, Th. 5); thus, they are themselves countable at infinity, which implies that X/G is paracompact (loc. cit.). Hence, there exists a locally finite covering \((\mathrm{U}_{\alpha}^{\prime})_{\alpha\in\mathrm{I}}\) of X/G by relatively compact open sets, and a covering \((\mathrm{V}_{\alpha}^{\prime})_{\alpha\in\mathrm{I}}\) such that \(\overline{V}_{\alpha}^{\prime}\subset U_{\alpha}^{\prime}\) for all \(\alpha\in I\) (General Topology, Chap. IX, §4, no. 3, Cor. 1 of Th. 3); taking the inverse image, we obtain two locally finite coverings \((\mathrm{U}_{\alpha})_{\alpha\in\mathrm{I}}\) and \((\mathrm{V}_{\alpha})_{\alpha\in\mathrm{I}}\) of X by relatively compact open sets stable under G, such that \(\overline{V}_{\alpha}\subset U_{\alpha}\) for all \(\alpha\in I\) .

For all \(\alpha \in \mathrm{I}\) , there exists a representation \(\rho_{\alpha}\) of G on a finite dimensional real vector space \(\mathrm{E}_{\alpha}\) and a morphism \(\varphi_{\alpha} \in \mathcal{C}^r (\mathrm{X};\mathrm{E}_{\alpha})\) , compatible with the operations of G, whose restriction to \(\mathrm{U}_{\alpha}\) is an embedding (Th. 1). For \(\alpha \in \mathrm{I}\) , let \(a_{\alpha}\) be a numerical function of class \(\mathrm{C}^r\) on X, equal to 1 on \(\mathrm{V}_{\alpha}\) and to 0 outside \(\mathrm{U}_{\alpha}\) (Differentiable and Analytic Manifolds, Results, 5.3.6). Put \(b_{\alpha}(x) = \int_{\mathrm{G}} a_{\alpha}(gx) dg\) for \(x \in \mathrm{X}\) . The function \(b_{\alpha}\) is of class \(\mathrm{C}^r\) , invariant under G (§6, no. 4, Cor. 2), equal to 1 on \(\mathrm{V}_{\alpha}\) and to 0 outside \(\mathrm{U}_{\alpha}\) . Give each \(\mathrm{E}_{\alpha}\) a Hilbert scalar product invariant under G (§1, no. 1), and R its canonical Hilbert structure; let E be the Hilbert sum of the family \((\mathrm{E}_{\alpha} \oplus \mathbf{R})_{\alpha \in \mathrm{I}}\) , and let \(\rho\) be the representation of G on E induced by the \(\rho_{\alpha}\) and the trivial action of G on R. For \(x \in \mathrm{X}\) , put \(\varphi(x) = (b_{\alpha}(x)\varphi_{\alpha}(x), b_{\alpha}(x))_{\alpha \in \mathrm{I}}\) . Then \(\varphi\) is a morphism of class \(\mathrm{C}^r\) from X to E, compatible with the operations of G; we verify as in the proof of Prop. 4 (no. 1) that \(\varphi\) is an embedding, which implies the corollary.

